\documentclass[12pt]{article}

% Layout.
\usepackage[top=1in, bottom=0.75in, left=1in, right=1in, headheight=1in, headsep=6pt]{geometry}

% Fonts.
\usepackage{mathptmx}
\usepackage[scaled=0.86]{helvet}
\renewcommand{\emph}[1]{\textsf{\textbf{#1}}}

% TiKZ.
\usepackage{tikz, pgfplots}
\usetikzlibrary{calc}

% Misc packages.
\usepackage{amsmath,amssymb,latexsym}
\usepackage{graphicx}
\usepackage{array}
\usepackage{xcolor}
\usepackage{multicol}
\usepackage{fancyhdr}
\usepackage{tikz}

%%% Force San Serif everywhere
\renewcommand{\familydefault}{\sfdefault}

\pagestyle{fancy}


% Misc.
\renewcommand{\d}{\displaystyle}
\newcommand{\ds}{\displaystyle}
\def\bc{\begin{center}}
\def\ec{\end{center}}
\def\be{\begin{enumerate}}
\def\ee{\end{enumerate}}

\newcommand{\ans}[1][1in]{\rule{#1}{.5pt}}


\lhead{Math F113X: Quiz 7}

\begin{document}

\ 

%\vspace{.25cm}

\noindent \textbf{Name:} \hrulefill \quad  score:\ans[1cm] \ / 10 \\


There are 10 points possible on this quiz. No aids (book, notes, etc.)
are permitted. You may use a non-programmable calculator.  {\bf Show
all work and supporting calculations for full credit. Explain how you
get your answers.}



\be

\item (2 points) \textbf{Decrypt} the message XUWNS LNXMJ WJ using an
  alphabetic Caesar (shift) cipher with shift 5 (A to F).
  % Spring is here

\vfill

\item (2 points)  \textbf{Encrypt} the message ``2 WEEKS LEFT'' using the substitution table below.\\

{
\begin{tabular}{|c|c|c|c|c|c|c|c|c|c|c|c|c|c|c|c|c|c|c|}
\hline
%&0&1&2&3&4&5&6&7&8&9&10&11&12&13&14&15&16&17\\ \hline \hline
plain &A&B&C&D&E&F&G&H&I&J&K&L&M&N&O&P&Q&R\\  \hline
cipher &L&K&J&I&P&O&N&M&T&S&Q&R&V&U&Y&X&W&Z \\
\hline
\end{tabular}

\begin{tabular}{|c|c|c|c|c|c|c|c|c|c|c|c|c|c|c|c|c|c|c|}
 \hline
%&18&19&20&21&22&23&24&25&26&27&28&29&30&31&32&33&34&35\\ \hline \hline
plain &S&T&U&V&W&X&Y&Z&0&1&2&3&4&5&6&7&8&9\\ \hline
cipher &9&8&7&6&5&4&3&2&1&0&D&C&B&A&H&G&F&E\\
%[12pt]
\hline
\end{tabular}
}

\vfill
\item (2 points)  \textbf{Encrypt} the message I LIKE CHEESE using a tabular transposition cipher with rows of length 4.
  

\vfill

\newpage

\item (2 points)   \textbf{Decrypt} the message LIKED ETESS BXRIZ using a tabular transposition cipher with rows of length 5.
%LETS RIDE BIKES

\vfill
\item (2 point) You are given a message to decrypt. You suspect that a
  substitution cipher may have been used to encrypt the message. Given
  this, you construct the following frequency table for the
  ciphertext message. What letter in plaintext might you guess is represented by
  ``X'' in ciphertext? Why?

  
	\begin{tabular}{c|c||c|c||c|c}
	letter&frequency &letter&frequency&letter&frequency\\
	\hline \hline
	A&4&J&12&S&0\\
	\hline
	B&13&K&20&T&21\\
	\hline
	C&14&L&9&U&5\\
	\hline
	D&9&M&1&V&15\\
	\hline
	E&0&N&0&W&1\\
	\hline
	F&0&O&6&X&38\\
	\hline
	G&8&P&8&Y&2\\
	\hline
	H&12&Q&5&Z&17\\
	\hline
	I&1&R&1&&\\
	\hline
	\end{tabular}

        \vspace{0.5cm}

        PLAINTEXT LETTER REPRESENTED BY ``X'':

        \vspace{0.5cm}

        EXPLANATION:

        \vspace{3cm}

        
\ee

\newpage

\makeatletter
\newcommand{\Letter}[1]{\@Alph{#1}}
\makeatother

\def\r{.6}
\begin{tikzpicture}
\draw[step=\r] (0,0) grid (28*\r,28*\r);
%\foreach \i in {1,2,...,26}{\draw (0,0) -- (\i*\r, 1);}
\foreach \i in {1,2,...,26}{\path let \n1 = {int(\i-1)} in (\i*\r+1/2*\r+1*\r, 28*\r-1/2*\r) node {\n1};}
\foreach \i in {1,2,...,26}{\path (\i*\r+1/2*\r +1*\r, 27*\r-1/2*\r) node {\textbf{\Letter{\i}}};}%
\foreach \i in {1,2,...,26}{\path (1/2*\r+1*\r, 26*\r-\i*\r+1/2*\r) node {\textbf{\Letter{\i}}};}%
\foreach \i in {1,2,...,26}{\path let \n1 = {int(\i-1)} in  (1/2*\r, 26*\r-\i*\r+1/2*\r) node {\n1};}%
\draw[line width = .75mm] (0,26*\r) -- (28*\r, 26*\r);
\draw[line width = .75mm] (\r+1*\r,0) -- (\r+1*\r, 28*\r);
\foreach \i in {1,2,...,26}
	\foreach \j in {1,2,...,26}{
		{\path let \n1 = {int(mod((\j-1)+(\i-1), 26)+1} in (\i*\r+1/2*\r+1*\r, 27*\r-\j*\r - 1/2*\r) node {\Letter{\n1}};
		}}
\end{tikzpicture}

\end{document}
%%% Local Variables:
%%% mode: LaTeX
%%% TeX-master: t
%%% End:
